\section{Capability Matching and Routing}

\subsection{Query Demand and Node Supply}

The routing problem can be described as a matching problem between query demand
and node supply. Let the set of supported capabilities be
\[
  C = \{\text{general}, \text{math}, \text{programming}, \text{writing},
  \text{summarization}, \text{research}, \text{planning}, \text{creative}\}.
\]

This set is deliberately compact, not a universal taxonomy of all LLM skills.
It covers common request families: information seeking, content processing or
summarization, procedural help, writing or rewriting, coding, quantitative
reasoning, and creative generation. These families are consistent with
user-need, instruction-task, and evaluation work
\citep{shelby2025tuna,wang2022supernaturalinstructions,liang2022helm}. The
small size also keeps DHT lookup and scoring simple.

A query is represented by a demand vector \(q \in [0,1]^{|C|}\), where
\(|C|\) is the number of supported capabilities. In normal operation, this
vector is produced by a capability classifier. In the implementation, it is
stored as a sparse map: omitted capabilities are treated as zero. The
classifier uses the user request and returns one to three scored capabilities
from the fixed set. For example, a simple math query may mostly need math,
while a prompt such as ``design a creative math puzzle and explain its
solution'' may need math, creative generation, and writing together.

Each node \(n\) also has capability scores \(s_n\). The value \(s_{n,c}\) means:
how good node \(n\) is for capability \(c\). The implemented prototype treats
these scores as part of the node profile; in this prototype, they are configured
or simulated rather than learned from a full model evaluation.

To compare a node with a query, the router computes a weighted average. For
each capability, it multiplies how much the query needs that capability by how
good the node is at it. Capabilities that matter more for the query therefore
count more in the final score:
\[
  Q(n, q) =
  \frac{\sum_{c \in C} q_c s_{n,c}}{\sum_{c \in C} q_c}.
\]

The division by \(\sum_{c \in C} q_c\) is a normalization step. The classifier
scores are treated as demand weights, not necessarily as percentages that
already sum to 1. Dividing by the total demand keeps \(Q(n,q)\) on the same
0--1 scale for simple and mixed queries. Without this normalization, a query
with three requested capabilities could receive a larger score than a
single-capability query only because more weights were added, not because the
node is actually a better match.

For example, if a query receives need scores of 0.9 for math and 0.6 for
writing, and a node scores 0.8 in math and 0.4 in writing, then
\[
  Q = \frac{0.9 \times 0.8 + 0.6 \times 0.4}{0.9 + 0.6} = 0.64.
\]
The node is therefore a 0.64 capability match for that query. The same formula
also works when the scores happen to sum to 1; in that case, the denominator is
1 and the expression reduces to the usual weighted sum.

This score is intentionally simple: interpretable, easy to inspect, and easy to
extend. It also bridges to richer profile-based routing systems. PPAI-style
query-agent matching, RouteProfile-style model profiles, and universal routing
can be viewed as richer alternatives for the same interface
\citep{wang2026ppai,xu2026routeprofile,jitkrittum2025uniroute}.

The demand vector has two roles. First, its strongest capabilities guide DHT
discovery: the router searches for peers advertising capabilities that look
useful for the request. Second, the complete vector is used for scoring. A
peer found through \texttt{programming}, for example, is still evaluated on the
full request if the query also needs writing or planning.

\subsection{Candidate Scoring}

The prototype does not rank peers only by raw capability. Routing work shows
that model selection trades off quality, cost, latency, and confidence
\citep{chen2023frugalgpt,ong2025routellm,moslem2026dynamicrouting}. Here, the
``best'' node is the node with the best routing score: capability quality is
combined with freshness, recent failures, and latency through a configurable
policy.

\[
  R(n, q) =
  \alpha Q(n, q)
  + \beta F(n)
  - \gamma E(n)
  - \delta L(n).
\]

Here \(F(n)\) is a freshness score, \(E(n)\) represents recent failures, and
\(L(n)\) represents latency. The weights define a policy. The default
\texttt{current} policy uses a strong capability weight with equal smaller
terms for freshness, failures, and latency:

\[
  R_{\text{current}}(n,q) =
  0.70Q(n,q) + 0.10F(n) - 0.10E(n) - 0.10L(n).
\]

Other implemented policies include \texttt{capability\_only},
\texttt{latency\_aware}, \texttt{reliability\_aware}, and \texttt{balanced}.
This makes the router a small experimental platform for studying matching
trade-offs.

The goal is therefore not simply to find the absolute best capability score for
every query. A node may have a strong model but still be a poor choice if it is
slow, stale, or often unavailable. Conversely, a suitable local node can score
well because it avoids network latency. This already helps the system avoid
pure quality maximization, while future policies can add explicit load and
request-spreading signals on the same scoring interface.

\begin{figure}[htbp]
\centering
\begin{tikzpicture}[
  signal/.style={draw=deepblue, rounded corners=2pt, thick, fill=softblue,
    align=center, minimum width=2.25cm, minimum height=0.62cm, font=\small},
  penalty/.style={draw=grayblack, rounded corners=2pt, thick, fill=graybg,
    align=center, minimum width=2.25cm, minimum height=0.62cm, font=\small},
  score/.style={draw=epflred, rounded corners=2pt, thick, fill=softred,
    align=center, minimum width=2.55cm, minimum height=0.72cm, font=\small},
  arrow/.style={-Latex, thick, color=grayblack}
]
\node[signal] (cap) at (0,1.55) {Capability\\quality \(Q\)};
\node[signal] (fresh) at (0,0.52) {Profile\\freshness \(F\)};
\node[penalty] (lat) at (0,-0.52) {Latency\\penalty \(L\)};
\node[penalty] (fail) at (0,-1.55) {Failure\\penalty \(E\)};
\node[score] (route) at (3.65,0) {Routing\\score \(R\)};
\node[signal] (select) at (7.1,0) {Selected node\\receives query\\and calls backend};

\draw[arrow] (cap.east) -- ([yshift=0.34cm]route.west);
\draw[arrow] (fresh.east) -- ([yshift=0.11cm]route.west);
\draw[arrow] (lat.east) -- ([yshift=-0.11cm]route.west);
\draw[arrow] (fail.east) -- ([yshift=-0.34cm]route.west);
\draw[arrow] (route.east) -- (select.west);
\end{tikzpicture}
\caption{The implemented routing score is transparent by design: capability is the main signal, while freshness, failures, and latency make operational health visible before the selected node receives the query.}
\end{figure}

\begin{table}[htbp]
\centering
\caption{Routing policies implemented by the prototype.}
\begin{tabularx}{\textwidth}{lrrrrX}
\toprule
\textbf{Policy} & \(\alpha\) & \(\beta\) & \(\gamma\) & \(\delta\) & \textbf{Intent} \\
\midrule
current & 0.70 & 0.10 & 0.10 & 0.10 & Prioritize capability while keeping freshness, failures, and latency visible. \\
capability\_only & 1.00 & 0.00 & 0.00 & 0.00 & Measure pure quality matching without health signals. \\
latency\_aware & 0.60 & 0.05 & 0.05 & 0.30 & Penalize slower peers more aggressively. \\
reliability\_aware & 0.60 & 0.05 & 0.30 & 0.05 & Penalize recent failures more aggressively. \\
balanced & 0.55 & 0.15 & 0.15 & 0.15 & Combine capability with equal freshness, failure, and latency signals. \\
\bottomrule
\end{tabularx}
\end{table}

\subsection{Routing Algorithm}

Chapter 4 described the architecture components. This subsection describes the
routing process in more detail: how those components are used by the
implemented router for one query.

The router follows a staged process. In the implementation, the routing service
selects a target, then the query service applies that decision by executing
locally or forwarding the query to the selected peer.

\begin{enumerate}
  \item Classify the user query into scored capability needs. This is the first
  LLM-backed step in a normal query: a classifier LLM, guided by a system
  prompt, maps the request to capability scores. For controlled evaluation, the
  required capabilities can also be injected directly. The implementation
  normalizes the result and keeps at most the top three capabilities.
  \item Query the DHT for providers of the relevant capabilities. The router
  searches capabilities with demand at least 0.3; if all scores are below this
  discovery threshold, it still searches the strongest capability.
  \item Fetch full profiles from the DHT providers and store or update them in
  the peer registry.
  \item Build the pre-recommendation candidate set from the local node and
  available registry profiles that match at least one requested capability. In a
  fresh run, these profiles mostly come from the current DHT lookup; over time,
  the registry can also contain profiles learned earlier. All candidates are
  scored again for the current query.
  \item If the best pre-recommendation candidate has a routing score of at least
  0.65, select it as the target.
  \item Otherwise, ask eligible remote candidates for recommendations. These
  recommendations expand the candidate set; recommended peers are then evaluated
  with the same routing score.
  \item If no candidate reaches 0.65, fall back to the best candidate seen. If
  no candidate was found at all and the original query did not already include
  \texttt{general}, retry discovery once with the general capability. If this
  retry also finds nothing, return \texttt{no\_suitable\_node}.
  \item Deliver the query to the selected target. If the target is the entry
  node itself, the query is executed locally. If the target is remote, the entry
  node opens a query protocol stream and sends the prompt, query identifier, and
  capability context to that peer.
  \item Generate and return the answer. This is the second LLM-backed step in a
  normal query: the selected node calls its bound answer backend, returns the
  answer to the entry node if needed, and the entry node returns it to the user.
\end{enumerate}

The 0.3 discovery threshold is used only to decide which capabilities should be
searched in the DHT. It avoids querying the DHT for very weak capability hints,
while the strongest capability is still searched if all scores are below 0.3.
The 0.65 routing threshold is a practical cutoff: high enough to avoid accepting
a weak pre-recommendation candidate immediately, but low enough to keep the
router from asking for recommendations too often. These values were chosen
empirically for this prototype and are exposed in the routing trace, together
with whether the selected peer came from the first candidate pool,
recommendations, or fallback selection.

Forwarding is intentionally not recursive. When a selected remote node receives
a forwarded query, it answers with its own backend instead of running the whole
routing process again. This avoids routing loops and unbounded delegation. If
forwarding fails, the entry node can exclude that peer and ask the router for
another target, but forwarding attempts are capped.

\subsection{Scope of the Current Router}

The implemented router is intentionally a single-executor router. For one user
request, it chooses one target node, sends the query there if needed, and returns
one answer. It does not yet split a task across several nodes, aggregate several
answers, use reputation, price, or privacy constraints, or optimize global
network load. Those extensions fit the same matching interface, but they belong
to the future-work layer rather than the current routing algorithm.
